\section{Methodology}
\label{sec:method}

In this section, we present the formal mathematical framework and architectural specifications of
\textbf{Task-Space Anchor Regularization (TSAR)}. We first describe the low-dimensional state
representation layer and task-space feature anchor extraction, followed by our dynamic linear
routing formulation, and conclude with the joint regularized calibration objective.

\subsection{Low-Dimensional State Representation}
To construct a computationally lightweight router, we avoid processing high-dimensional
intermediate visual activations directly. Let $z(x)_b \in \mathbb{R}^D$ represent the feature
representation mathematically simulated to mimic the globally pooled activations of a Vision
Transformer (ViT-Tiny) backbone for calibration sample $b$, where $D=192$. 

Using a subset of the calibration split $X_{cal}$, we compute a frozen, unsupervised Principal
Component Analysis (PCA) projection matrix $P \in \mathbb{R}^{D \times d}$ to project features into
a compact, low-dimensional routing subspace of dimension $d = K = 4$. Specifically, the
high-dimensional visual features $z(x)_b$ are mean-centered over the calibration split prior to
computing the singular value decomposition (SVD) of the sample covariance matrix, with the
projection matrix $P$ comprising the top $d$ right-singular vectors. 

While computing PCA on a tiny dataset ($B_{cal} = 64$ samples total) would normally introduce high
sampling noise and parameter variance, the pre-trained backbone representations exhibit high
inter-task separability. Because task-specific representations lie in highly distinct coordinate
spaces, the top eigenvectors of the covariance matrix represent exceptionally stable multi-task axes
that are easily captured from a handful of samples. To completely eliminate any potential risk of
sampling noise or meta-overfitting on the projection space, practitioners can also: (i) compute the
PCA matrix $P$ over a larger task-agnostic unlabeled pre-training set, or (ii) employ a completely
data-independent random Gaussian projection (Johnson-Lindenstrauss projection), which mathematically
guarantees pairwise distance preservation within a bounded distortion factor without requiring any
calibration data (as analyzed in Section \ref{sec:sandbox_extrap}). Indeed, under extreme data
scarcity (e.g., $B_{cal} \le 32$), our empirical ablation in Appendix~\ref{sec:projection_ablation}
shows that data-independent Random Gaussian projection consistently and substantially outperforms
unsupervised PCA (+5.26\% Joint Mean accuracy at $B_{cal}=16$ while cutting seed variance by more
than half) because it is entirely immune to local calibration sampling noise. We therefore recommend
Random Gaussian projection as the preferred default scheme for ultra-sparse low-data constraints.

To stabilize optimization and ensure scale invariance across disparate visual tasks, we normalize
the projected coordinate onto the unit sphere. The low-dimensional normalized routing state
$\psi(x)_b \in \mathbb{R}^d$ is formulated as:
\begin{equation}
\psi(x)_b = \frac{z(x)_b P}{\|z(x)_b P\|_2 + \epsilon}
\label{eq:state_proj}
\end{equation}
where $\epsilon = 10^{-8}$ is a small positive constant to prevent division by zero. Note that
while the PCA projection matrix $P$ is computed on mean-centered calibration features to align with
standard Singular Value Decomposition (SVD) conventions, the forward projection in
Equation~\ref{eq:state_proj} is applied directly to the raw feature vector $z(x)_b$ without explicit
test-time mean subtraction. Let $\mu_z \in \mathbb{R}^D$ represent the global mean of the
calibration features, such that the true centered and projected coordinates are formulated as
$(z(x)_b - \mu_z)P$. Applying Equation~\ref{eq:state_proj} to uncentered features is mathematically
equivalent to introducing a translation vector $\mu_P = \mu_z P$ inside both the numerator and the
$L_2$ norm divisor:
\begin{equation}
\psi(x)_b = \frac{(z(x)_b - \mu_z) P + \mu_P}{\|(z(x)_b - \mu_z) P + \mu_P\|_2 + \epsilon}
\label{eq:uncentered_proj_approx}
\end{equation}
Because the translation vector $\mu_P$ resides inside the norm divisor, the normalization scales
each sample non-linearly depending on the magnitude of the centered projected vector
$(z(x)_b - \mu_z) P$, introducing a sample-dependent non-linear distortion. Mathematically, this is
not strictly equivalent to a constant linear shift that can be perfectly absorbed by the downstream
bias parameters $B_{l,k}$ (or $B_k$). However, under our representation sandbox, the task manifolds
are highly separated on the unit sphere, meaning the centered projected vector's norm
$\|(z(x)_b - \mu_z) P\|_2$ is highly concentrated. Under this regime, the non-linear distortion is
exceptionally small and behaves as a harmless, highly localized scaling perturbation. This
approximation completely eliminates the operational overhead of storing, tracking, and subtracting
global feature means during real-time edge system deployment, while preserving outstanding
empirical routing stability.

\subsection{Task Feature Anchors}
A major challenge of low-data calibration is that the coordinate space of the projected features
$\psi(x)_b$ exhibits high variance under sample scarcity, making the routing optimization prone to
overfitting local noise. Our key observation is that the centroids of pre-trained expert feature
representations represent stable, task-specific coordinates (or anchors) that can guide
optimization.

Let $X_{cal, k}$ denote the partition of the calibration set $X_{cal}$ containing samples belonging
to task $k \in \{1, \dots, K\}$. We compute the task-space feature anchor $\bar{\psi}_k \in
\mathbb{R}^d$ by averaging the low-dimensional routing states of all samples belonging to task $k$:
\begin{equation}
\bar{\psi}_k = \frac{1}{|X_{cal, k}|} \sum_{b \in X_{cal, k}} \psi(x)_b
\label{eq:anchor_mean}
\end{equation}
These centroids $\bar{\psi}_k$ are computed offline once prior to router calibration, requiring
zero additional overhead during the optimization loops or inference.

\subsection{Routing Coefficient Dynamics}
For a given input sample $b$, we instantiate a compact, single-layer global linear router ($L=1$) as our recommended default architecture. The dynamic, sample-specific merging coefficient $\alpha_{k, b}$ is computed as:
\begin{equation}
\alpha_{k, b}^{TSAR} = \langle \psi(x)_b, W_{k} \rangle + B_{k}
\label{eq:router_coeffs_l1}
\end{equation}
where $W_{k} \in \mathbb{R}^d$ and $B_{k} \in \mathbb{R}$ represent the trainable routing weights and biases, respectively. For a task count of $K=4$ and projection dimension $d=4$, the complete routing parameter footprint is only 20 parameters, ensuring near-zero computational overhead during both optimization and inference deployment.

While the single-layer global router is the optimal and recommended choice for classification head model merging (where coefficient collapse renders layer-wise routers redundant, as proven below), we also formulate the generalized multi-layer router ($L > 1$) to maintain mathematical generality and support layer-wise deep model merging (where intermediate layer weights are merged using layer-specific routing coefficients). In this generalized formulation, for layer $l \in \{1, \dots, L\}$, the dynamic, sample-specific merging coefficient $\alpha_{k, b}(l)$ is computed as:
\begin{equation}
\alpha_{k, b}^{TSAR}(l) = \langle \psi(x)_b, W_{l, k} \rangle + B_{l, k}
\label{eq:router_coeffs}
\end{equation}
where $W_{l, k} \in \mathbb{R}^d$ and $B_{l, k} \in \mathbb{R}$ represent the trainable layer-wise weights and biases.

Unlike softmax-based routers, our formulation is linear and unconstrained, which prevents the
zero-sum competitive bottleneck identified in prior literature \cite{trial5_submission5}, preserving
classical representation capacity. 

In our controlled representation sandbox, the specialized expert weights $W_k$ represent the joint
effective classification parameters of the model. To combine these experts dynamically under our generalized multi-layer formulation, we pool the layer-wise dynamic router coefficients across the $L=14$ layer groups. Under batch deployment, we average these coefficients across a batch of size $B$ and across all layer groups $L$:
\begin{equation}
\bar{\alpha}_k = \frac{1}{L} \sum_{l=1}^L \left( \frac{1}{B} \sum_{b=1}^B \alpha_{k, b}(l) \right)
\label{eq:batch_coeff}
\end{equation}
The merged model parameter matrix $W_{merged}$ and bias vector $B_{merged}$ are then dynamically
constructed as a weighted linear combination of the pre-trained expert parameter matrices $W_k$ and
bias vectors $B^{expert}_k$:
\begin{equation}
W_{merged} = \sum_{k=1}^K \bar{\alpha}_k W_k, \quad B_{merged} = \sum_{k=1}^K \bar{\alpha}_k
B^{expert}_k
\label{eq:merging_sum}
\end{equation}

\paragraph{Mathematical Equivalence to Output-Level Logit Ensembling:}
We must highlight a fundamental mathematical property of our sandbox. Because the specialized
expert models are linear classifiers applied directly on top of a frozen feature space, the output
logits for an input feature representation $z$ under the merged parameters are:
\begin{align}
\text{logits} &= W_{\text{merged}} z + B_{\text{merged}} \nonumber \\
&= \left( \sum_{k=1}^K \bar{\alpha}_k W_k \right) z + \left( \sum_{k=1}^K \bar{\alpha}_k
B^{\text{expert}}_k \right) \nonumber \\
&= \sum_{k=1}^K \bar{\alpha}_k \left( W_k z + B^{\text{expert}}_k \right)
\label{eq:logit_equivalence}
\end{align}
This derivation reveals that in our linear sandbox, dynamic parameter-level model merging is
mathematically identical to taking a weighted linear combination of the output logits of each
individual expert model (dynamic output-level logit ensembling). 

We embrace this equivalence as a major analytical strength. It makes the parameter-merging process
mathematically tractable and transparent, allowing us to isolate and study routing mechanics,
gradient cross-talk, and stream audits independently of architectural confounders. 

However, we stress that for multi-layer deep neural networks with non-linear activation functions
(e.g., ReLUs, GeLUs, or self-attention blocks), parameter-level model merging and output-level logit
ensembling are physically distinct and diverge fundamentally:
\begin{align}
\sigma\big( W_{\text{merged}, 2} &\sigma( W_{\text{merged}, 1} x ) \big) \nonumber \\
&\neq \sum_{k=1}^K \bar{\alpha}_k \sigma\big( W_{k, 2} \sigma( W_{k, 1} x ) \big)
\label{eq:deep_merging_divergence}
\end{align}
In deep networks, merging weights of intermediate layers alters the representation flow
non-linearly, which cannot be replicated by linear combinations at the output level. Thus, while our
linear sandbox serves as a controlled proxy to study calibration optimization, the actual physical
challenges of model merging in deep networks are much more complex, and we include output-level
logit ensembling as a key baseline in our evaluation to empirically verify this mathematical
equivalence.

Crucially, we must analyze the structural implications of this pooling formulation. Mathematically,
substituting the sample-specific linear routing equation (Equation \ref{eq:router_coeffs}) into our
batch-average pooling equation (Equation \ref{eq:batch_coeff}) yields:
\begin{align}
\bar{\alpha}_k &= \frac{1}{L} \sum_{l=1}^L \left( \frac{1}{B} \sum_{b=1}^B \left( \langle \psi(x)_b, W_{l, k} \rangle + B_{l, k} \right) \right) \nonumber \\
&= \frac{1}{B} \sum_{b=1}^B \left( \langle \psi(x)_b, \frac{1}{L} \sum_{l=1}^L W_{l, k} \rangle +
\frac{1}{L} \sum_{l=1}^L B_{l, k} \right)
\label{eq:collapse_proof}
\end{align}
Let us define the globally collapsed parameters as the layer-wise averages:
\begin{equation}
W_{\text{global}, k} = \frac{1}{L} \sum_{l=1}^L W_{l, k}, \quad B_{\text{global}, k} = \frac{1}{L}
\sum_{l=1}^L B_{l, k}
\label{eq:collapsed_params}
\end{equation}
where $W_{\text{global}, k} \in \mathbb{R}^d$ and $B_{\text{global}, k} \in \mathbb{R}$. We can
then rewrite the final pooled dynamic coefficient as:
\begin{equation}
\bar{\alpha}_k = \frac{1}{B} \sum_{b=1}^B \left( \langle \psi(x)_b, W_{\text{global}, k} \rangle + B_{\text{global}, k} \right)
\label{eq:collapsed_coeff}
\end{equation}
This derivation formally exposes the phenomenon of \textbf{layer-averaging collapse}. At inference
deployment, our layer-wise router collapses to a single global low-dimensional router, rendering the
$L=14$ layer-wise weight divisions mathematically redundant in terms of representation capacity. 

However, from an optimization perspective, training this $L$-layer over-parameterized routing model
on a tiny 64-sample calibration split introduces key mathematical advantages over direct
single-layer optimization:
\begin{enumerate}
  \item \textbf{Gradient Damping and Variance Reduction (Gradient Bagging):} Under severe low-data
calibration, optimization gradients are exceptionally noisy, making a 20-parameter single-layer
router highly prone to bad local minima or early divergence. In our $L$-layer over-parameterized
formulation, backpropagation gradients with respect to each individual layer weight $W_{l, k}$ are
damped by the scaling factor $1/L$ via the chain rule:
  \begin{equation}
  \frac{\partial \mathcal{L}_{CE}}{\partial W_{l, k}} = \frac{1}{L} \frac{\partial
\mathcal{L}_{CE}}{\partial \bar{\alpha}_k} \psi(x)_b
  \label{eq:damped_grad}
  \end{equation}
  This $1/L$ gradient damping, combined with randomized independent parameter initializations,
smooths the optimization trajectory. Individual layer-wise routers act as distinct "routing heads"
exploring different local basins of the loss landscape, and their final average $W_{global, k}$ acts
as a robust ensemble (akin to bagging across layers), delivering substantially lower gradient
variance and superior seed stability.
  \item \textbf{Distributed Geometric Springs:} Under TSAR, the regularization loss (Equation
\ref{eq:tsar_penalty}) acts as a distributed system of geometric springs, pulling each of the $L$
layers independently toward the stable task anchor $\bar{\psi}_k$. This multi-layer regularized
constraint stabilizes optimization path trajectories, preventing high-frequency parameter
oscillations and enabling the ensemble average parameters $W_{global, k}$ to converge robustly onto
the true task coordinates.
\end{enumerate}

\subsection{TSAR Calibration Loss Function}
To calibrate the routing weights $W_{l, k}$ and biases $B_{l, k}$ under extreme data scarcity, we
augment the standard multi-task cross-entropy loss over the calibration split with a spatial
distance penalty that pulls the routing weights $W_{l, k}$ toward their respective task feature
anchors $\bar{\psi}_k$. The total joint optimization objective is defined as:
\begin{equation}
\mathcal{L}_{total} = \mathcal{L}_{CE} + \mathcal{L}_{WD} + \mathcal{L}_{TSAR}
\label{eq:total_loss}
\end{equation}
where:
\begin{equation}
\mathcal{L}_{WD} = \lambda_{wd} \sum_{l=1}^L \sum_{k=1}^K \left( \|W_{l, k}\|_2^2 + B_{l, k}^2 \right)
\label{eq:l2_wd}
\end{equation}
is the standard $L_2$ weight decay regularization, and:
\begin{equation}
\mathcal{L}_{TSAR} = \lambda_{anchor} \sum_{l=1}^L \sum_{k=1}^K \| W_{l, k} - \bar{\psi}_k \|_2^2
\label{eq:tsar_penalty}
\end{equation}
is our proposed Task-Space Anchor Regularization penalty, scaled by the coefficient
$\lambda_{anchor} \ge 0$. Calibration is performed for $100$ epochs using the AdamW optimizer with a
learning rate of $1.0 \times 10^{-2}$ and weight decay coefficient $\lambda_{wd} = 10^{-3}$ (see
Appendix~\ref{sec:reproducibility} for complete hyperparameter specifications).

By forcing the router parameters to stay aligned with the task feature anchors, TSAR acts as a
geometric spatial prior. It ensures that the routing weights do not scale excessively or drift into
arbitrary noise-aligned subspaces during low-data optimization, thereby maintaining stable,
task-aligned routing gradients.

\subsection{Multi-Task Gradient Balancing via PCGrad}
While the TSAR spatial penalty prevents parameter drift, dynamic router calibration under
multi-task setups is still subject to severe gradient cross-talk. Specifically, when task
difficulties differ significantly (e.g., MNIST is simple, whereas SVHN is highly noisy and complex),
the total calibration gradient $\nabla_{W} \mathcal{L}_{total}$ is heavily dominated by the
gradients of the hardest tasks. This gradient imbalance can drive the routing parameters of simple
tasks away from their stable anchors to minimize hard-task loss, leading to severe optimization
instability.

To resolve this multi-task gradient cross-talk, we implement \emph{Projecting Conflicting
Gradients} (PCGrad) \cite{yu2020gradient} during calibration. Instead of backpropagating the
averaged cross-entropy loss directly, we compute the individual task losses $\mathcal{L}^{(k)}$
separately on the calibration samples. For each task $k \in \{1, \dots, K\}$, we compute its
corresponding gradient vector $g_k = \nabla_{W} \left( \mathcal{L}^{(k)} + \frac{1}{K}
(\mathcal{L}_{WD} + \mathcal{L}_{TSAR}) \right)$. 

We then identify conflicts between task gradients and project them to preserve mutual optimization
direction. If two gradients $g_i$ and $g_j$ conflict (indicated by a negative cosine similarity,
i.e., $g_i \cdot g_j < 0$), we project $g_i$ onto the normal plane of $g_j$:
\begin{equation}
g_i \leftarrow g_i - \frac{g_i \cdot g_j}{\|g_j\|_2^2} g_j
\label{eq:pcgrad_proj}
\end{equation}
This projection is applied iteratively across tasks in a randomized order. The final balanced
gradient $g = \sum_k g_k$ is then applied to update the router parameters. By explicitly projecting
conflicting gradients, PCGrad preserves cross-task inhibition—preventing dominant tasks from
corrupting easily aligned routing parameters—while allowing all tasks to be mutually optimized under
extreme data constraints. While PCGrad introduces an $O(K)$ computational complexity scaling
bottleneck during training, we present a comprehensive scalability analysis and propose three highly
practical mitigations (Task Grouping, Selective Projection, and Stochastic Task Sampling) in
Appendix~\ref{sec:pcgrad_complexity} to keep the main text self-contained and ensure real-world
scalability. Specifically, when $K \ge 20$, task grouping reduces the backward passes to $O(G)$
(where $G \ll K$ is the number of semantic clusters), while stochastic task sampling maintains a
constant $O(1)$ scaling cost per step, rendering PCGrad computationally viable for massive
multi-task systems.

\subsection{Architecture and Complexity Specifications}
Our recommended default architecture is a compact single-layer global router ($L=1$) with a projection dimension $d=4$. Since the projection dimension matches the number of tasks $K=4$, the complete routing parameter footprint is:
\begin{itemize}
  \item \textbf{Weights ($W$):} Shape $(K \times d) = (4 \times 4) = 16$ parameters.
  \item \textbf{Biases ($B$):} Shape $(K) = 4$ parameters.
  \item \textbf{Total Router Parameters:} 20 trainable parameters.
\end{itemize}

For research purposes and to support deep layer-wise model merging of intermediate weights (where different weights are merged at each layer group and layer-wise routing is active), we also support the generalized $L=14$ layer-wise router:
\begin{itemize}
  \item \textbf{Weights ($W$):} Shape $(L \times K \times d) = (14 \times 4 \times 4) = 224$ parameters.
  \item \textbf{Biases ($B$):} Shape $(L \times K) = (14 \times 4) = 56$ parameters.
  \item \textbf{Total Router Parameters:} 280 trainable parameters.
\end{itemize}

Both parameter footprints are ultra-lightweight, representing less than $0.005\%$ of the backbone model's parameter size, ensuring near-zero computational overhead during backpropagation and deployment inference.